\documentclass{article}
%DIF LATEXDIFF DIFFERENCE FILE
%DIF DEL /Users/ntsoi/Downloads/STARS_ From Spatiotemporal Dynamics to Social Representations in Human-Robot Interaction (Version 9889)/main_new.tex   Tue Jun  2 13:59:16 2026
%DIF ADD /Users/ntsoi/Downloads/main_3_new.tex                                                                                                         Tue Jun  2 09:22:00 2026

% Datasets
% CMU TBD Dataset
% UT CODa
% UT Scand
% Waymo
% Nu Scenes
% *new* Collected via LiDAR and Cameras on UT Campus
% SEAN Together

\usepackage{xcolor}
\usepackage{amsmath}
\usepackage{amssymb}
\usepackage{amsfonts}
\usepackage{subcaption}
\usepackage{graphicx}
\usepackage{booktabs}
\usepackage{tabularx}
\usepackage{multirow}
\usepackage{xspace}
\usepackage{enumitem}
\usepackage{svg}
\usepackage{xr}
% This macro tells latexmkrc that main.tex depends on sup.tex
\makeatletter
\newcommand*{\addFileDependency}[1]{%
  \typeout{(#1)}%
  \@addtofilelist{#1}%
  \IfFileExists{#1}{}{\typeout{No file #1.}}%
}
\makeatother

\newcommand*{\myexternaldocument}[1]{%
  \externaldocument[sup-]{#1}%  Adds 'sup-' prefix to avoid label name clashes
  \addFileDependency{#1.tex}%
  \addFileDependency{#1.aux}%
}
\myexternaldocument{sup}

\newcommand{\todo}[1]{\textcolor{red}{\textbf{TODO: #1}}}


\usepackage{corl_2026} % Use this for the initial submission.
% \usepackage[final]{corl_2026} % Uncomment for the camera-ready ``final'' version.
% \usepackage[preprint]{corl_2026} % Uncomment for pre-prints (e.g., arxiv); This is like ``final'', but will remove the CORL footnote.

\newcommand{\method}{STARS\xspace}
\newcommand{\methodfull}{\textbf{S}patio\textbf{T}emporal \textbf{A}utoencoded \textbf{R}epresentations for \textbf{S}ocial-interaction\xspace}
\newcommand{\pretraining}{Spatiotemporal Graph Autoencoding\xspace}


\newcommand{\jb}[1]{\textcolor{red}{JB: #1]}}
\newcommand{\nathan}[1]{\textcolor{green}{[Nathan: #1]}}
\newcommand{\mike}[1]{\textcolor{blue}{[Mike: #1]}}
\newcommand{\tejas}[1]{\textcolor{purple}{[Tejas: #1]}}
\newcommand{\rishab}[1]{\textcolor{cyan}{[Rishab: #1]}}


\title{STARS: From Spatiotemporal Dynamics to Social Representations in Human-Robot Interaction}


% The \author macro works with any number of authors. There are two
% commands used to separate the names and addresses of multiple
% authors: \And and \AND.
%
% Using \And between authors leaves it to LaTeX to determine where to
% break the lines. Using \AND forces a line break at that point. So,
% if LaTeX puts 3 of 4 authors names on the first line, and the last
% on the second line, try using \AND instead of \And before the third
% author name.

% NOTE: authors will be visible only in the camera-ready and preprint versions (i.e., when using the option 'final' or 'preprint'). 
% 	For the initial submission the authors will be anonymized.

\author{
  Jane E.~Doe\\
  Department of Electrical Engineering and Computer Sciences\\
  University of California Berkeley 
  United States\\
  \texttt{janedoe@berkeley.edu} \\
}
%DIF PREAMBLE EXTENSION ADDED BY LATEXDIFF
%DIF UNDERLINE PREAMBLE %DIF PREAMBLE
\RequirePackage[normalem]{ulem} %DIF PREAMBLE
\RequirePackage{color}\definecolor{RED}{rgb}{1,0,0}\definecolor{BLUE}{rgb}{0,0,1} %DIF PREAMBLE
\providecommand{\DIFadd}[1]{{\protect\color{blue}\uwave{#1}}} %DIF PREAMBLE
\providecommand{\DIFdel}[1]{{\protect\color{red}\sout{#1}}} %DIF PREAMBLE
%DIF SAFE PREAMBLE %DIF PREAMBLE
\providecommand{\DIFaddbegin}{} %DIF PREAMBLE
\providecommand{\DIFaddend}{} %DIF PREAMBLE
\providecommand{\DIFdelbegin}{} %DIF PREAMBLE
\providecommand{\DIFdelend}{} %DIF PREAMBLE
\providecommand{\DIFmodbegin}{} %DIF PREAMBLE
\providecommand{\DIFmodend}{} %DIF PREAMBLE
%DIF FLOATSAFE PREAMBLE %DIF PREAMBLE
\providecommand{\DIFaddFL}[1]{\DIFadd{#1}} %DIF PREAMBLE
\providecommand{\DIFdelFL}[1]{\DIFdel{#1}} %DIF PREAMBLE
\providecommand{\DIFaddbeginFL}{} %DIF PREAMBLE
\providecommand{\DIFaddendFL}{} %DIF PREAMBLE
\providecommand{\DIFdelbeginFL}{} %DIF PREAMBLE
\providecommand{\DIFdelendFL}{} %DIF PREAMBLE
\providecommand{\DIFscaledelfig}{0.5}
%DIF HIGHLIGHTGRAPHICS PREAMBLE %DIF PREAMBLE
\RequirePackage{settobox} %DIF PREAMBLE
\RequirePackage{letltxmacro} %DIF PREAMBLE
\newsavebox{\DIFdelgraphicsbox} %DIF PREAMBLE
\newlength{\DIFdelgraphicswidth} %DIF PREAMBLE
\newlength{\DIFdelgraphicsheight} %DIF PREAMBLE
% store original definition of \includegraphics %DIF PREAMBLE
\LetLtxMacro{\DIFOincludegraphics}{\includegraphics} %DIF PREAMBLE
\providecommand{\DIFaddincludegraphics}[2][]{{\color{blue}\fbox{\DIFOincludegraphics[#1]{#2}}}} %DIF PREAMBLE
\providecommand{\DIFdelincludegraphics}[2][]{% %DIF PREAMBLE
\sbox{\DIFdelgraphicsbox}{\DIFOincludegraphics[#1]{#2}}% %DIF PREAMBLE
\settoboxwidth{\DIFdelgraphicswidth}{\DIFdelgraphicsbox} %DIF PREAMBLE
\settoboxtotalheight{\DIFdelgraphicsheight}{\DIFdelgraphicsbox} %DIF PREAMBLE
\scalebox{\DIFscaledelfig}{% %DIF PREAMBLE
\parbox[b]{\DIFdelgraphicswidth}{\usebox{\DIFdelgraphicsbox}\\[-\baselineskip] \rule{\DIFdelgraphicswidth}{0em}}\llap{\resizebox{\DIFdelgraphicswidth}{\DIFdelgraphicsheight}{% %DIF PREAMBLE
\setlength{\unitlength}{\DIFdelgraphicswidth}% %DIF PREAMBLE
\begin{picture}(1,1)% %DIF PREAMBLE
\thicklines\linethickness{2pt} %DIF PREAMBLE
{\color[rgb]{1,0,0}\put(0,0){\framebox(1,1){}}}% %DIF PREAMBLE
{\color[rgb]{1,0,0}\put(0,0){\line( 1,1){1}}}% %DIF PREAMBLE
{\color[rgb]{1,0,0}\put(0,1){\line(1,-1){1}}}% %DIF PREAMBLE
\end{picture}% %DIF PREAMBLE
}\hspace*{3pt}}} %DIF PREAMBLE
} %DIF PREAMBLE
\LetLtxMacro{\DIFOaddbegin}{\DIFaddbegin} %DIF PREAMBLE
\LetLtxMacro{\DIFOaddend}{\DIFaddend} %DIF PREAMBLE
\LetLtxMacro{\DIFOdelbegin}{\DIFdelbegin} %DIF PREAMBLE
\LetLtxMacro{\DIFOdelend}{\DIFdelend} %DIF PREAMBLE
\DeclareRobustCommand{\DIFaddbegin}{\DIFOaddbegin \let\includegraphics\DIFaddincludegraphics} %DIF PREAMBLE
\DeclareRobustCommand{\DIFaddend}{\DIFOaddend \let\includegraphics\DIFOincludegraphics} %DIF PREAMBLE
\DeclareRobustCommand{\DIFdelbegin}{\DIFOdelbegin \let\includegraphics\DIFdelincludegraphics} %DIF PREAMBLE
\DeclareRobustCommand{\DIFdelend}{\DIFOaddend \let\includegraphics\DIFOincludegraphics} %DIF PREAMBLE
\LetLtxMacro{\DIFOaddbeginFL}{\DIFaddbeginFL} %DIF PREAMBLE
\LetLtxMacro{\DIFOaddendFL}{\DIFaddendFL} %DIF PREAMBLE
\LetLtxMacro{\DIFOdelbeginFL}{\DIFdelbeginFL} %DIF PREAMBLE
\LetLtxMacro{\DIFOdelendFL}{\DIFdelendFL} %DIF PREAMBLE
\DeclareRobustCommand{\DIFaddbeginFL}{\DIFOaddbeginFL \let\includegraphics\DIFaddincludegraphics} %DIF PREAMBLE
\DeclareRobustCommand{\DIFaddendFL}{\DIFOaddendFL \let\includegraphics\DIFOincludegraphics} %DIF PREAMBLE
\DeclareRobustCommand{\DIFdelbeginFL}{\DIFOdelbeginFL \let\includegraphics\DIFdelincludegraphics} %DIF PREAMBLE
\DeclareRobustCommand{\DIFdelendFL}{\DIFOaddendFL \let\includegraphics\DIFOincludegraphics} %DIF PREAMBLE
%DIF AMSMATHULEM PREAMBLE %DIF PREAMBLE
\makeatletter %DIF PREAMBLE
\let\sout@orig\sout %DIF PREAMBLE
\renewcommand{\sout}[1]{\ifmmode\text{\sout@orig{\ensuremath{#1}}}\else\sout@orig{#1}\fi} %DIF PREAMBLE
\makeatother %DIF PREAMBLE
%DIF COLORLISTINGS PREAMBLE %DIF PREAMBLE
\RequirePackage{listings} %DIF PREAMBLE
\RequirePackage{color} %DIF PREAMBLE
\lstdefinelanguage{DIFcode}{ %DIF PREAMBLE
%DIF DIFCODE_UNDERLINE %DIF PREAMBLE
  moredelim=[il][\color{red}\sout]{\%DIF\ <\ }, %DIF PREAMBLE
  moredelim=[il][\color{blue}\uwave]{\%DIF\ >\ } %DIF PREAMBLE
} %DIF PREAMBLE
\lstdefinestyle{DIFverbatimstyle}{ %DIF PREAMBLE
	language=DIFcode, %DIF PREAMBLE
	basicstyle=\ttfamily, %DIF PREAMBLE
	columns=fullflexible, %DIF PREAMBLE
	keepspaces=true %DIF PREAMBLE
} %DIF PREAMBLE
\lstnewenvironment{DIFverbatim}[1][]{\lstset{style=DIFverbatimstyle}}{} %DIF PREAMBLE
\lstnewenvironment{DIFverbatim*}[1][]{\lstset{style=DIFverbatimstyle,showspaces=true}}{} %DIF PREAMBLE
\lstset{extendedchars=\true,inputencoding=utf8}

%DIF END PREAMBLE EXTENSION ADDED BY LATEXDIFF

\begin{document}
\maketitle

\begin{abstract}
Fielding socially competent robots requires joint reasoning over spatial and social information. While current models excel at either spatial reasoning (e.g., modeling physical dynamics) or social reasoning (e.g., parsing text-based interactions), few natively unify both. Moreover, data scarcity in Human-Robot Interaction (HRI) makes training large models from scratch challenging. To bridge this gap, we introduce \textbf{STARS}: \textbf{S}patio\textbf{T}emporal \textbf{A}utoencoded \textbf{R}epresentations for \textbf{S}ocial-interaction, a method for learning \emph{event-level relational representations} from HRI datasets. Rather than training task-specific models, STARS models fixed-duration interaction scenarios as spatiotemporal graphs. A self-supervised message-passing graph neural network autoencoder then compresses these physical dynamics into a compact latent space for each node and edge. This injects strong relational inductive bias to naturally capture the underlying social situation. We evaluate STARS' representations on the SEAN Together dataset, which provides VR-based navigation data annotated with subjective human perceptions, and the SocialNav-SUB benchmark for social scene understanding. Linear probing demonstrates that STARS achieves highly competitive performance on perception prediction and pedestrian action classification tasks. Specifically, STARS exhibits excellent data efficiency, achieving an $F_1$-Score of 0.635 on predicting perceived competence using just 1\% of labeled data,
whereas baselines require up to ten times as much \DIFaddbegin \DIFadd{labeled data to achieve comparable performance}\DIFaddend . Furthermore, qualitative analysis of the learned latent space reveals that STARS naturally encodes high-level social semantics, demonstrating cohesive clustering of distinct human action intents without explicit task supervision.
\end{abstract}

\keywords{Social Robotics, Human-Robot Interaction, Representation Learning, Graph Neural Networks} 

\section{Introduction}
Deploying socially competent robots in human spaces requires autonomous robots to reason jointly using spatial knowledge to plan collision-free paths and social knowledge to perceive context, predict behaviors, and adhere to social norms. A robot navigating through a crowded hallway, for example, must plan collision-free motion in the physical environment while also anticipating pedestrian behavior and reacting in a socially compliant manner~\cite{mavrogiannis2023core, francis2025principles}.
Joint reasoning over spatial and social information is central to Human-Robot Interaction (HRI), yet existing approaches tend to address only one or the other.
Foundation spatial models parse LiDAR and visual observations into geometric representations, but lack the capacity to understand social dynamics~\cite{zhou2025social, lisondra2026embodied}.
At the same time, some autoregressive models conditioned on text (LLMs) or text and vision (VLMs) exhibit strong social perception~\cite{zhou2025social}, but struggle to process continuous spatial dynamics~\cite{ramakrishnan2024does, munje2025socialnavsub} and precise temporal information~\cite{gao2025vision}.
Bridging this gap by training large models from scratch is challenging due to the scarcity of grounded HRI data.
Collecting such data requires time and cost intensive human-subject studies, and yields datasets that are small, fragmented across sensor modalities, and highly prone to overfitting when applied to large, unstructured models.

Although most HRI datasets are small and fragmented across settings and tasks, the premise of this paper is that their recurring social interaction structure~\cite{de2015makes} can be used to learn a unified representation that generalizes across diverse downstream social tasks that require reasoning about agents, interactions, and shared context. This interaction structure is naturally relational, as agents move through shared spaces and influence one another over time. We hypothesize that graph-structured representations are well-suited to learning from small, fragmented HRI datasets because they can encode agents as nodes and interactions as edges, injecting strong \textit{relational inductive bias}~\cite{battaglia2018relational}. This inductive bias enables message-passing graph neural networks~\cite{gilmer2020message} to learn interaction patterns efficiently in limited-data regimes.

We introduce \textbf{\method}: \methodfull, a method for unsupervised learning of \emph{event-level relational representations}.
In contrast to the autoregressive prediction of world models, \method learns a representation over a short time horizon to capture relevant contextual information.
We define an \emph{interaction scenario} as a fixed-duration sequence of events where agents share a physical space.
We encode these interaction scenarios into graphs, where directed edges encode sequences of pairwise spatial features over time.

Though our interaction scenario graphs contain only spatiotemporal information, our target downstream tasks are inherently social.
By optimizing a self-supervised reconstruction objective over the fixed-duration interaction scenario, \method compresses physical dynamics into a compact latent space that captures the implicit \textit{social situation}~\cite{tsoi2022sean}.
We demonstrate that the learned representation allows a simple linear probe to extract social meaning from physical motion, unlocking data-efficient generalization without requiring task-specific representation learning.

Our contributions are threefold: 
\vspace{-0.5em}
\begin{enumerate}[nosep, leftmargin=*]
    \item We introduce \method for unsupervised representation learning to compress HRI datasets of various modalities into a unified latent representation containing features useful for diverse downstream tasks in HRI while requiring minimal task-specific labeled data.
    \item We demonstrate that incorporating relational inductive bias improves data efficiency. By comparing \method to unstructured baselines in low-data regimes, we show that explicitly restricting the model architecture to represent entities and their interactions enables robust generalization with a fraction of the labeled data.
    \item We establish that \method' self-supervised representations naturally encode social semantics. By analyzing the latent space topology, we illustrate the model's inherent ability to distinguish complex social behaviors, such as action intents, without explicit supervision.
\end{enumerate}


\section{Related Work}
\label{sec:relatedwork}

\paragraph{Spatial and social reasoning.}

Evaluating multi-agent interactions has historically relied on pedestrian trajectory datasets~\cite{pellegrini2009you, lerner2007crowds, robicquet2016learning} and egocentric robot observations~\cite{martin2019jrdb, karnan2022socially}.
To model these complex shared environments, Graph Neural Networks (GNNs) have emerged as the standard architecture, as they naturally mirror the topological structure of human crowds~\cite{mohamed2020social, kosaraju2019social}.
Recent heterogeneous and attention-based spatiotemporal GNNs excel at weighting relational dependencies between diverse entities to extract continuous kinematic interactions~\cite{li2025pedestrian, li2025unified}.
Current learning efforts tend to focus on trajectory forecasting, which aims to minimize spatial displacement error, and therefore learn to reflect future spatial configurations, rather than reusable social semantics.
Recent works have laid the groundwork to push beyond spatial forecasting.
\citet{zhang2023predicting} provide a dataset and baseline for predicting human perceptions of robot performance and \citet{munje2025socialnavsub} established a benchmark for VLM understanding of social interactions in navigation.
Still, existing trajectory-focused models struggle to generalize to \textit{social} tasks because they do not explicitly model the normative behaviors and interpersonal associations that define social interactions~\cite{thompson2025social, tsoi2022sean}.

\paragraph{Foundation models and the HRI data bottleneck.}
Foundation and world models demonstrate impressive predictive capabilities for embodied decision-making~\cite{bommasani2021opportunities, hafner2023mastering, bruce2024genie, parkerholder2024genie2}, including strong zero-shot transfer in Vision-Language-Action (VLA) controllers~\cite{o2024open, zitkovich2023rt, kim2024openvla}.
However, these models rely on massive, unified datasets.
In contrast, grounded HRI data is scarce, geographically constrained, and highly fragmented across varying sensor modalities, survey instruments, and tasks~\cite{duncan2024survey, ondras2022human, bu2024ssup, heinisch2024physiological,heo2024diverse, thompson2021conversational, yang2021dataset, shrestha2024natsgd}.
Applying foundation-scale learning to HRI is therefore hindered by data scarcity, leading unstructured models to overfit or fail to generalize across diverse social scenarios~\cite{lisondra2026embodied}.

\paragraph{Relational inductive bias.}
To learn efficiently in low-data regimes, recent approaches have leveraged structured representations.
Relational inductive bias~\cite{battaglia2018relational} constrains neural architectures to represent entities as nodes and interactions as edges, encouraging networks to deduce underlying physical dynamics without relying on hand-engineered features~\cite{ferreira2019learning}.
While this bias has driven advancements in physical object manipulation, its application to extracting social context from interaction remains underexplored~\cite{sanchez2020learning, malik2023relational}.
At the same time, generative self-supervised graph representation learning has shown that semantically meaningful embeddings can be extracted without human annotation by reconstructing masked features or topological structures~\cite{kipf2016variational, liu2025network, falqueto2025human}.


\section{Method}
\label{sec:method}

\DIFaddbegin \DIFadd{In this section, we detail the }\method \DIFadd{framework.
We first formalize the problem of extracting representations from multi-modal HRI datasets, outline our graph construction process, and describe the unsupervised graph autoencoder used to learn the latent social semantics.
}

\DIFaddend \subsection{Problem Formulation}
Let $\mathcal{D} = \{\mathcal{D}_1, \mathcal{D}_2, \dots, \mathcal{D}_K\}$ denote a collection of HRI datasets (Fig.~\ref{fig:ds_sns}A).
For the purposes of this work, we consider an HRI dataset to be any dataset containing observations of humans and robots, or humans in shared physical spaces relevant to robot interaction.
Each dataset $\mathcal{D}_k$ contains a set of interaction scenarios $\{x_i^{(k)}\}_{i=1}^{N_k}$.
An interaction scenario $x$ is a fixed-duration observation of $T$ timesteps containing multiple agents sharing a physical space.
Depending on the dataset, $x$ \DIFdelbegin \DIFdel{may include }\DIFdelend \DIFaddbegin \DIFadd{includes features such as }\DIFaddend trajectories, poses, velocities, robot state, \DIFdelbegin \DIFdel{local maps, or other sensor-derived features}\DIFdelend \DIFaddbegin \DIFadd{or local maps}\DIFaddend .
Datasets may differ in sensor modality, sampling frequency, agent types, and available annotations.
\DIFaddbegin \DIFadd{To handle these diverse inputs, all raw sensor modalities are projected into a common hidden dimension prior to message passing.
}\DIFaddend 

The representation-learning problem maps an interaction scenario $x \in \mathcal{X}$ to a compact representation $\mathbf{z} \in \mathcal{Z}$ that supports downstream HRI tasks.
For a downstream task $\mathcal{T}_m$, a subset of scenarios is paired with labels $\{(x_i, y_i^{(m)})\}$,
where $y_i^{(m)} \in \mathcal{Y}_m$ \DIFdelbegin \DIFdel{may represent }\DIFdelend \DIFaddbegin \DIFadd{represents the task-specific annotation, which in our evaluations include }\DIFaddend subjective human perceptions \DIFdelbegin \DIFdel{, social actions, group membership, or other task-specific annotations}\DIFdelend \DIFaddbegin \DIFadd{of robot behavior and discrete pedestrian action categories}\DIFaddend .
A representation is considered better if it yields stronger held-out performance on $\mathcal{T}_m$ under the same labeled-data budget and downstream adaptation protocol.
Thus, the goal is to obtain compact and reusable representations that support strong downstream performance when labeled examples for each task are scarce.

\begin{figure}[p!tb]
    \centering
    \includegraphics[width=\textwidth]{figures/method.pdf}
    \caption{
    \textbf{\method} three stages:
    (\textbf{A}) construct spatiotemporal graphs from multi-modal Human-Robot Interaction datasets,
    (\textbf{B}) utilize an unsupervised message-passing graph autoencoder to compress physical dynamics into latent representations, and
    (\textbf{C}) adapt frozen latents to downstream tasks using lightweight task-specific prediction heads.
    See Sec. \ref{sec:method} for details.
    }
    \label{fig:method}
\end{figure}

\subsection{\method Overview} 
Motivated by a need to learn from diverse HRI datasets containing unlabeled spatiotemporal interactions while preserving the relational structure needed for downstream social tasks, \method is 
designed to disentangle spatiotemporal dynamics via autoencoding to latent vectors $\mathbf{z} = f_\phi(x)$ that capture reusable social context~\cite{tsoi2022sean}.

As shown in Figure~\ref{fig:method}, \method \DIFdelbegin \DIFdel{enables the learning of representations for social tasks }\DIFdelend \DIFaddbegin \DIFadd{is a two-stage framework for learning representations }\DIFaddend from spatiotemporal data\DIFdelbegin \DIFdel{in three stages}\DIFdelend .
First, \method extracts features from multi-modal HRI datasets (examples shown in Figure~\ref{fig:ds_sns}) to construct unified spatiotemporal graphs.
While \method applies generally to any spatiotemporal relational data, in this work we focus on the application area of social navigation.
Explicitly modeling human and robot agents as nodes and their physical interactions as edges imposes a strong relational inductive bias.
Second, a pretraining step compresses physical dynamics into a compact latent space via a self-supervised message-passing graph autoencoder.
Self-supervised reconstruction reduces overfitting to scarce labels, enabling higher-capacity encoders (Supp. Mat. Sec.~\ref{sup-sec:supervised_heterograph_ablation}).
\DIFdelbegin \DIFdel{Third, }%DIFDELCMD < \method %%%
\DIFdel{adapts the pretrained latent representation to tasksby learning a lightweight}\DIFdelend \DIFaddbegin \DIFadd{To evaluate the utility of these representations for social tasks, we subsequently apply a lightweight, }\DIFaddend task-specific \DIFdelbegin \DIFdel{prediction head}\DIFdelend \DIFaddbegin \DIFadd{linear prediction head to the frozen latents}\DIFaddend .
We instantiate \method using both homogeneous and heterogeneous graph formulations to study how explicit entity and relation types affect \method' learned representations.

\subsection{Spatiotemporal Relational Graph Construction}

Given a fixed-duration interaction scenario of $T$ timesteps, \method represents the scenario as a directed spatiotemporal graph $G=(V, E)$. Each node $v_i \in V$ corresponds to an agent, and each directed edge $e_{ij} \in E$ represents the ordered pairwise interaction from source agent $i$ to target agent $j$ over the fixed-duration interaction scenario. For each edge $e_{ij}$, we define a temporal edge feature sequence $\mathbf{e}_{ij} = [\mathbf{e}_{ij,1}, \dots, \mathbf{e}_{ij,T}] \in \mathbb{R}^{T \times d_e}$, where $\mathbf{e}_{ij,t}$ contains pairwise spatial and motion features at timestep $t$, such as relative position, relative heading, and relative velocity.
Our construction makes two design choices.
First, our graphs are \textit{edge-feature dense}, capturing temporal relationships primarily through edge rather than node attributes.
Second, rather than using distance-based thresholds, we construct \textit{complete graphs} to capture interactions between all observed agents.

\textbf{The homogeneous graph} represents nodes as generic ``agents'' and all edges ``interactions.''
Robot and human nodes retain distinct input attributes, but the graph does not include type embeddings.

\textbf{The heterogeneous graph}
treats all nodes by their entity types.
Edges are also typed, given the connected nodes.
The node set is partitioned as $V = V_R \cup V_H$, where $V_R$ contains robot nodes and $V_H$ contains human nodes.
Robot nodes may include robot-specific attributes, such as encoded occupancy grids or trajectory goals when available, while human nodes are associated with a learnable semantic embedding representing the entity type.
The edge set captures typed relations such as robot-human, human-robot, and human-human interactions, with relation types provided separately through edge-type embeddings.

\subsection{Unsupervised Graph Autoencoder}

The core of \method is a graph encoder implemented as a Message-Passing Graph Neural Network (MPGNN) that maps the spatiotemporal graph into node-level latent distributions.

\textbf{Feature Preprocessing:}
Node attributes are projected into a shared hidden space via an MLP.
A shared Transformer edge encoder $\text{Enc}_{\text{edge}}$ with learned positional embeddings maps each temporal edge sequence $\mathbf{e}_{ij}$ to a fixed-dimensional relation embedding~\cite{vaswani2017attention,devlin2019bert}.

\textbf{Message Passing:}
The graph encoder performs message passing over the complete interaction graph.
In \method' \textit{heterogeneous} graph formulation, message passing uses a HEAT-style update~\cite{mo2021heat}.
Let $\mathbf{h}_j^{(k)}$ and $\mathbf{h}_i^{(k)}$ denote the source and destination node embeddings at layer $k$, and $\mathbf{t}_{ij}$ denote a learned embedding for the specific edge type (e.g., human-robot vs. human-human).
%
We construct a joint interaction context:
\begin{equation}
\mathbf{u}_{ij}^{(k)} = [\mathbf{h}_j^{(k)} \| \mathbf{h}_i^{(k)} \| \text{Enc}_{\text{edge}}(\mathbf{e}_{ij}) \| \mathbf{t}_{ij}]
\end{equation}
%
The representation $\mathbf{u}_{ij}^{(k)}$ includes the full context of the interaction.
During message passing, the network must determine both the content of the information being transmitted and the relative importance of that information.
Therefore, we use $\mathbf{u}_{ij}^{(k)}$ to independently parameterize both the raw attention score $s_{ij}^{(k)}$ and the message content $\mathbf{m}_{ij}^{(k)}$ via two learnable functions:
\begin{equation}
s_{ij}^{(k)} = f_{\mathrm{attn}}(\mathbf{u}_{ij}^{(k)}), \qquad \mathbf{m}_{ij}^{(k)} = f_{\mathrm{msg}}(\mathbf{u}_{ij}^{(k)})
\end{equation}
Attention weights are normalized via a destination-wise softmax $\alpha_{ij}^{(k)}$, and the aggregated message is used to update the node states via a residual LayerNorm step; see Sec.~\ref{sup-sec:message-passing} in the Supplementary Material for details.
Typed-edge attention enables the model to learn which distant interactions matter without needing separate message-passing modules per relation.
The \textit{homogeneous} graph formulation omits edge-type embeddings $\mathbf{t}_{ij}$ from $\mathbf{u}_{ij}^{(k)}$, using only the source node embedding, destination node embedding, and encoded temporal edge features $\text{Enc}_{\text{edge}}(\mathbf{e}_{ij})$ in the update.

\textbf{Latent Encoding:} After $K$ message-passing layers, shared linear projection heads parameterize a diagonal Gaussian latent distribution for each node:
\begin{equation}
q_\phi(\mathbf{z}_i \mid G) = \mathcal{N}\!\left(\boldsymbol{\mu}_i,\; \mathrm{diag}(\boldsymbol{\sigma}_i^2)\right), \qquad \boldsymbol{\mu}_i, \log \boldsymbol{\sigma}_i^2 = f_\mu(\mathbf{h}_i^{(K)}),\; f_\sigma(\mathbf{h}_i^{(K)})
\end{equation}
Latent samples are drawn via the reparameterization trick, $\mathbf{z}_i = \boldsymbol{\mu}_i + \boldsymbol{\sigma}_i \odot \boldsymbol{\epsilon}$.
This per-node latent formulation preserves the asymmetry of pairwise interactions, as edge reconstruction conditions on both endpoint latent nodes rather than a single pooled graph representation.

\textbf{Decoding \& Objective:}
Unlike autoregressive models, we employ an unsupervised autoencoding objective over the entire fixed-duration interaction scenario.
Specifically, for temporal edge reconstruction, the decoder employs an MLP that takes the concatenated source and destination latents $[\mathbf{z}_i \| \mathbf{z}_j]$ as input. 
This MLP projects the joint latent representation into a hidden dimension and outputs a single flat vector (e.g., $\in \mathbb{R}^{320}$) for each edge. 
This sequence-level output is then reshaped into the explicit temporal sequence format $T \times d_e$ (e.g., $40 \times 8$) to yield the reconstructed edge features $\mathbf{\hat{e}}_{ij}$.

For node-level reconstruction, we dynamically restrict the objective to nodes containing active, non-trivial input attributes.
Reconstructing static, zero-initialized features (e.g. human nodes in SEAN or all nodes in SNS) would introduce a strong bias toward trivial constant targets, diluting the gradient signal.
Let $V_{\text{feat}} = \{ v_i \in V \mid \mathbf{x}_i \neq \mathbf{0} \}$ be the subset of nodes containing active attributes.
In our datasets, this corresponds to the robot node features in SEAN (e.g., occupancy grids and goals).
For human nodes and SNS graphs, we omit node reconstruction.
Crucially, the absence of human node reconstruction does not exclude human nodes from the learned representations or lead to a training mismatch because edge features $\mathbf{e}_{ij}$ represent continuous relative pairwise interactions (e.g., relative positions, headings, and velocities over time) between all agents, and the temporal edge decoder must reconstruct $\mathbf{\hat{e}}_{ij}$ from the concatenated endpoint latents $[\mathbf{z}_i \| \mathbf{z}_j]$.
The model is forced to encode rich spatiotemporal social semantics and joint context into the human node latents to successfully reconstruct these edges.

The model is trained by minimizing a modified Evidence Lower Bound (ELBO):
\begin{equation}
\label{eq:elbo}
\mathcal{L}_{\text{pre}} = \sum_{v_i \in V_{\text{feat}}} \lambda_{v} \text{MSE}(\mathbf{x}_i, \mathbf{\hat{x}}_i) + \sum_{e_{ij} \in E} \lambda_{ij} \text{MSE}(\mathbf{e}_{ij}, \mathbf{\hat{e}}_{ij}) + \beta \sum_{v_i \in V} D_{\text{KL}} \left( q_\phi(\mathbf{z}_i \mid G) \parallel \mathcal{N}(0, I) \right)
\end{equation}
where $\lambda_{v}$ and $\lambda_{ij}$ denote weighting factors for node and edge reconstruction, allowing the network to learn across imbalanced datasets\DIFaddbegin \DIFadd{, for example, }\DIFaddend where human-human interactions are much more common than human-robot interactions.

\subsection{Task-Agnostic Linear Probing}

To evaluate the quality of our learned latent representations and adapt them to downstream tasks with limited labeled data, we apply linear probes~\cite{alain2016understanding}.
Because \method learns compressed node-level latent representations $\mathbf{z}_i$, it provides flexibility in selecting and aggregating these node latents to represent different structural components (such as edges or subgraphs) relevant to a given downstream task.

After freezing the pretrained encoder weights $\phi$, \method selects and aggregates the necessary latents based on a given downstream task.
For instance, a target-conditioned task like predicting human perception might concatenate the robot latent $\mathbf{z}_r$, a target human latent $\mathbf{z}_{h, \text{target}}$, and a mean-pooled representation of bystander latents $\overline{\mathbf{z}}_{h, \text{bystander}}$.
Conversely, a pairwise task such as group detection can represent an edge relationship by concatenating the node latents of the two endpoints (e.g., $[\mathbf{z}_i \parallel \mathbf{z}_j]$).
A linear probe with weights $W$ is then optimized using supervised pairs $(G, y)$, where $y$ is the task label:
%
\begin{equation}
\mathcal{L}_{\text{task}} = \text{LossFn} \left( y, W \mathbf{z}_{\text{probe}} + b \right)
\end{equation}
%
This lightweight adaptation step keeps the representation fixed, so downstream performance reflects the information encoded in $\mathbf{z}_{\text{probe}}$ rather than task-specific representation learning.

% Source is in figures/datasets.kra (a https://krita.org/en/ app file)
\begin{figure*}[t]
  \centering
  \includegraphics[width=\linewidth]{figures/ds-sns-legend-no-cpt.png}
  \caption{\textbf{(A) Dataset Overview.} 
  We utilize two human-robot interaction datasets: \textbf{SEAN-T}~\cite{zhang2023predicting} for evaluating subjective human ratings of robot navigation and \textbf{SNS}~\cite{munje2025socialnavsub} for downstream relational action classification (e.g., \textit{avoid}, \textit{follow}). 
  \textbf{(B) Latent Space Analysis (RQ3).}
  2D Projections of the SNS Human Node Action Latent Space for the \textit{avoid} and \textit{follow} actions.
  We visualize the 128-dimensional representations of nodes extracted from the \method-trained 
  homogeneous GNN-VAE encoder using linear \textbf{PCA} and non-linear \textbf{t-SNE}.
  Each point represents a node in the latent graph embedded, color-coded by action label.
  See Sec.~\ref{sec:latent_space_semantics} for details.}
  \label{fig:ds_sns}
\end{figure*}

\section{Experiments}
\label{sec:experiments}

We evaluate the effectiveness of \method to learn social semantics from physical dynamics.
Specifically, we demonstrate that \method achieves data-efficient downstream performance, and through analysis of the latent space topology, we show that the model can distinguish social behaviors, such as action intents, without explicit supervision during pretraining.
Our research questions are:

\begin{enumerate}[nosep, label={\textbf{RQ\arabic*:}}, leftmargin=*]
    \item \textbf{Downstream Performance.} How well does \method work on downstream HRI tasks compared to unstructured baselines when adapting pretrained representations?
    \item \textbf{Data Efficiency.} For a given downstream task, what is the impact of labeled data availability (from 1\% to 100\%) when using pretrained \method representations?
    \item \textbf{Latent Space Semantics.} Do \method' self-supervised representations naturally encode social semantics (e.g., action intents and group dynamics) without explicit supervision?
\end{enumerate}


\paragraph{Datasets and Tasks}
We evaluate \method across two distinct datasets to assess unsupervised pretraining and downstream task adaptation.
Our pretraining pool consists of unlabeled, multi-agent spatiotemporal data from \textbf{SEAN Together (SEAN-T)}~\cite{zhang2023predicting}, which provides 8-second immersive virtual-reality observations of robot navigation, and \textbf{SocialNav-SUB (SNS)}~\cite{munje2025socialnavsub}, a structured social reasoning benchmark containing pedestrian-robot interactions.
Examples of the interaction scenarios are shown in Figure~\ref{fig:ds_sns} and further details are provided in the Supp. Mat. Sec.~\ref{sup-sec:datasets}.
%
To evaluate whether our autoencoder successfully captures reusable social semantics, we freeze the pretrained representations and train lightweight linear probes on two distinct downstream tasks: 
%
1) \textbf{Perception Prediction (SEAN-T):} Predicting binary subjective human ratings of a robot's behavior across three distinct metrics: perceived competence, clarity of intention, and surprise~\cite{zhang2023predicting}. This task evaluates the model's ability to extract subjective social impressions directly from physical motion.
2) \textbf{Action Classification (SNS):} Predicting discrete, categorical relational actions taken by pedestrians in dense crowd scenes with up to 30 people (e.g., \emph{avoiding}, \emph{following}, \emph{yielding})~\cite{munje2025socialnavsub}. This task tests the model's capacity to explicitly classify interaction intents in complex multi-agent environments.

\paragraph{Training and Evaluation}
All models were trained and evaluated controlling for model capacity by keeping parameter counts comparable, utilized fast early stopping~\cite{prechelt2002early} to ensure convergence while preventing overfitting, utilized extensive hyperparameter search for all models, and were repeated 10 times to account for the variability in random weight initialization.
Complete training protocols are detailed in Supp. Mat. Sec.~\ref{sup-sec:training_eval_protocol}.


\subsection{RQ1 \& RQ2: Downstream Performance and Data Efficiency}
\label{sec:downstream_efficiency}

\paragraph{Experiment:}

To evaluate downstream performance (RQ1) and data efficiency (RQ2), we pretrain \method on the combined, unlabeled training splits of the datasets.
After freezing the encoder, we extract the relevant elements of the latent graph representation (details in the Supp. Mat. Sec.~\ref{sup-sec:training_eval_protocol}) and train the linear probe using progressively larger fractions of the labels from the training split.
We compare our approach against several unstructured baselines: a fully connected, feed-forward Multi-Layer Perceptron (MLP), an Autoencoder to ablate the graph structure while retaining data compression, and a Random Forest to provide a non-deep-learning baseline.
Note that the Random Forest is omitted for SNS due to the dataset's variable-sized, high-dimensional features.
We report $F_1$-Scores for \textit{competence}, \textit{surprise}, and \textit{intention} for SEAN-T and macro $F_1$-Score over \textit{avoid}, \textit{follow}, and \textit{not consider} for SNS, with per-class results reported in the Supp. Mat. Sec. \ref{sup-sec:extended_main_results}.


\paragraph{Results:}

Table~\ref{tbl:r1} shows the heterogeneous (\method-He) and homogeneous (\method-Ho) variants consistently outperform unstructured baselines on SEAN-T.
While both exhibit competitive results, \method-Ho yields the highest overall scores in most SEAN-T categories. 
However, \method-He demonstrates notable early efficiency in the Surprise prediction task and outperforms all models on the SNS action classification task, both in macro $F_1$ and across the individual SNS action classes reported in Supp. Mat. Sec.~\ref{sup-sec:extended_main_results}.
On SNS, unstructured baselines fail to extract relational signal, plateauing near a macro $F_1$-Score of $0.3$, regardless of how much labeled data is provided.
In contrast, \method-He rapidly acquires task-relevant information, achieving $0.362$ with just 1\% of the data and climbing to $0.503$ at 100\%, likely because explicit entity and edge typing helps it better distinguish robot–human interactions in dense crowd scenes.
Despite these task-specific variations, \method requires much less labeled data than baseline methods to achieve high performance.
The relational inductive bias inherent in \method allows the model to rapidly generalize even when restricted to 5\% or 10\% of the labeled downstream data. 
This advantage is especially notable under extreme low-data regimes; for instance, on the Competence task, \method-Ho achieves an $F_1$-Score of $0.612$ using only 1\% of the labeled data, whereas unstructured baselines like the Random Forest~\cite{tsoi2022sean} or MLP require roughly five to ten times as much data to reach comparable performance. 
Finally, across nearly all tasks and data regimes, a linear probe trained on \method's frozen representations outperforms unstructured baselines.
The one exception is the SEAN-T \textit{surprise} task at 10\% labeled data ($0.552 \pm 0.17$ Autoencoder vs. $0.548 \pm 0.20$ \method-Ho), though the difference falls well within a standard deviation, which is not unexpected given the naturally high variance of predicting subjective human surprise.

%%% SEAN-T & SNS RESULTS
\begin{table}[tb!p]
\centering
\caption{
Downstream Performance (RQ1) and Data Efficiency (RQ2) on SEAN-T and SNS.
Results showcase $F_1$-Scores ($\mu \pm \sigma$) across progressively scaled fractions of labeled data.
We compare unstructured baseline methods against our proposed \method variants (\textbf{\method-He}: Heterogeneous, \textbf{\method-Ho}: Homogeneous).
To rigorously evaluate both representation quality and data efficiency, \method utilizes a frozen pretrained encoder paired with a lightweight linear probe.} 
\label{tbl:r1}
\setlength{\tabcolsep}{6pt}
\resizebox{\textwidth}{!}{
\begin{tabular}{ll *{6}{c}} 
\toprule
 & \textbf{Method} & \textbf{1\%} & \textbf{5\%} & \textbf{10\%} & \textbf{25\%} & \textbf{50\%} & \textbf{100\%} \\
\midrule

% --- Competence Section ---
\multirow{5}{*}{\rotatebox[origin=c]{90}{\begin{tabular}{@{}c@{}}\textbf{SEAN-T} \\ \textbf{Competence}\end{tabular}}} 
& MLP   & $0.450 \pm 0.22$ & $0.546 \pm 0.19$ & $0.664 \pm 0.15$ & $0.761 \pm 0.03$ & $0.778 \pm 0.02$ & $0.797 \pm 0.02$ \\
& Autoencoder    & $0.528 \pm 0.18$ & $0.677 \pm 0.05$ & $0.710 \pm 0.04$ & $0.739 \pm 0.04$ & $0.776 \pm 0.02$ & $0.779 \pm 0.02$ \\
& Random Forest  & $0.557 \pm 0.20$ & $0.641 \pm 0.12$ & $0.680 \pm 0.10$ & $0.704 \pm 0.09$ & $0.715 \pm 0.09$ & $0.718 \pm 0.08$ \\
\cmidrule(lr){2-8}
& \textbf{\method-He (Ours)}& $\mathbf{0.635 \pm 0.10}$ & $0.673 \pm 0.07$ & $0.726 \pm 0.04$ & $0.762 \pm 0.03$ & $0.781 \pm 0.03$ & $0.794 \pm 0.02$ \\
& \textbf{\method-Ho (Ours)}  & $0.612 \pm 0.14$ & $\mathbf{0.702 \pm 0.08}$ & $\mathbf{0.738 \pm 0.05}$ & $\mathbf{0.765 \pm 0.04}$ & $\mathbf{0.787 \pm 0.03}$ & $\mathbf{0.799 \pm 0.02}$ \\
\midrule

% --- Surprise Section ---
\multirow{5}{*}{\rotatebox[origin=c]{90}{\begin{tabular}{@{}c@{}}\textbf{SEAN-T} \\ \textbf{Surprise}\end{tabular}}} 
& MLP   & $0.217 \pm 0.22$ & $0.250 \pm 0.14$ & $0.423 \pm 0.27$ & $0.450 \pm 0.21$ & $0.630 \pm 0.23$ & $0.720 \pm 0.09$ \\
& Autoencoder    & $0.372 \pm 0.21$ & $0.468 \pm 0.18$ & $\mathbf{0.552 \pm 0.17}$ & $0.558 \pm 0.12$ & $0.624 \pm 0.11$ & $0.721 \pm 0.05$ \\
& Random Forest  & $0.291 \pm 0.24$ & $0.341 \pm 0.19$ & $0.373 \pm 0.20$ & $0.444 \pm 0.16$ & $0.470 \pm 0.15$ & $0.496 \pm 0.15$ \\
\cmidrule(lr){2-8}
& \textbf{\method-He (Ours)}& $\mathbf{0.455 \pm 0.17}$ & $0.481 \pm 0.17$ & $0.526 \pm 0.15$ & $0.603 \pm 0.12$ & $0.655 \pm 0.08$ & $0.709 \pm 0.04$ \\
& \textbf{\method-Ho (Ours)}  & $0.440 \pm 0.23$ & $\mathbf{0.541 \pm 0.21}$ & $0.548 \pm 0.20$ & $\mathbf{0.610 \pm 0.16}$ & $\mathbf{0.713 \pm 0.06}$ & $\mathbf{0.741 \pm 0.03}$ \\
\midrule

% --- Intention Section ---
\multirow{5}{*}{\rotatebox[origin=c]{90}{\begin{tabular}{@{}c@{}}\textbf{SEAN-T} \\ \textbf{Intention}\end{tabular}}} 
& MLP   & $0.479 \pm 0.20$ & $0.533 \pm 0.13$ & $0.606 \pm 0.06$ & $0.659 \pm 0.03$ & $0.693 \pm 0.03$ & $0.703 \pm 0.01$ \\
& Autoencoder    & $0.546 \pm 0.08$ & $0.562 \pm 0.09$ & $0.635 \pm 0.03$ & $0.653 \pm 0.03$ & $0.694 \pm 0.02$ & $0.686 \pm 0.02$ \\
& Random Forest  & $0.592 \pm 0.18$ & $0.665 \pm 0.10$ & $0.685 \pm 0.09$ & $0.692 \pm 0.09$ & $0.700 \pm 0.08$ & $0.698 \pm 0.09$ \\
\cmidrule(lr){2-8}
& \textbf{\method-He (Ours)}& $0.576 \pm 0.08$ & $0.619 \pm 0.05$ & $0.671 \pm 0.04$ & $0.694 \pm 0.03$ & $0.705 \pm 0.03$ & $0.721 \pm 0.03$ \\
& \textbf{\method-Ho (Ours)}  & $\mathbf{0.603 \pm 0.10}$ & $\mathbf{0.671 \pm 0.07}$ & $\mathbf{0.701 \pm 0.06}$ & $\mathbf{0.715 \pm 0.04}$ & $\mathbf{0.723 \pm 0.03}$ & $\mathbf{0.734 \pm 0.02}$ \\
\midrule

% --- Macro F1 Section (SNS) ---
\multirow{5}{*}{\rotatebox[origin=c]{90}{\begin{tabular}{@{}c@{}}\textbf{SNS} \\ \textbf{Macro $F_1$}\end{tabular}}} 
& MLP    & $0.294 \pm 0.00$ & $0.297 \pm 0.01$ & $0.300 \pm 0.01$ & $0.306 \pm 0.02$ & $0.295 \pm 0.01$ & $0.293 \pm 0.00$ \\
& Autoencoder     & $0.297 \pm 0.01$ & $0.294 \pm 0.00$ & $0.294 \pm 0.00$ & $0.294 \pm 0.00$ & $0.294 \pm 0.00$ & $0.294 \pm 0.00$ \\
% & Random Forest  & --- & --- & --- & --- & --- & --- \\
\cmidrule(lr){2-8}
& \textbf{\method-He (Ours)} & $\mathbf{0.362 \pm 0.058}$ & $\mathbf{0.447 \pm 0.051}$ & $\mathbf{0.470 \pm 0.044}$ & $\mathbf{0.493 \pm 0.032}$ & $\mathbf{0.501 \pm 0.028}$ & $\mathbf{0.503 \pm 0.022}$ \\
& \textbf{\method-Ho (Ours)} & $0.344 \pm 0.05$ & $0.431 \pm 0.045$ & $0.453 \pm 0.035$ & $0.468 \pm 0.03$ & $0.471 \pm 0.024$ & $0.475 \pm 0.02$ \\
\bottomrule
\end{tabular}
}
\end{table}







\subsection{RQ3: Latent Space Semantics}
\label{sec:latent_space_semantics}

\paragraph{Experiment:}
To qualitatively evaluate whether the learned representations capture high-level social semantics without supervision (RQ3), we analyze the topology of the latent spaces.
Using the pretrained latent representation from the homogeneous graph encoder reported in Tbl. \ref{tbl:r1}, we extract frozen representations from the 128-dimensional node-level embeddings for SNS.
To focus on semantic navigation intents, we conduct our experiments on the dominant action labels in the SocialNav-Sub (SNS) dataset, \textit{avoid} and \textit{follow}.
We apply both Principal Component Analysis (PCA) for linear projection and t-Distributed Stochastic Neighbor Embedding (t-SNE) for non-linear dimensionality reduction to map the high-dimensional representations to a 2D space, color-coded by downstream class labels.

\paragraph{Results:}
Latent space projections shown in Figure~\ref{fig:ds_sns} reveal structural self-organization corresponding to downstream human navigation semantics. 
The linear PCA projection (left) displays a continuous, topologically smooth gradient across the state space, indicating that linear relationships and global dynamics are well preserved. 
In contrast, the non-linear t-SNE projection (right) reveals cohesive, well-separated clusters for action classes.
The semantic groupings show the homogeneous graph encoder's ability to distinguish complex human intent dynamics without relying on any downstream labels during pretraining.



\section{Limitations and Future Work}
\label{sec:limitations}
%\paragraph{Limitations and Future Work.}
While \method generalizes well, its autoencoder and message passing inherently smooth high-frequency pairwise geometries.
Therefore, these latents excel at node-level tasks like intent prediction, but are less effective on edge-centric tasks (e.g., group detection) that require exact relational geometry.
Additionally, standardizing diverse datasets into unified graphs can discard dataset-specific nuances.
Future work will explore more flexible graph construction pipelines and the use of intermediate, pre-message-passing temporal edge embeddings to bypass the autoencoder bottleneck, preserving the precise geometric information necessary for fine-grained relational tasks.

\section{Conclusion}
\label{sec:conclusion}
In this work, we introduced \method to overcome the HRI data bottleneck by bridging continuous physical dynamics and high-level social reasoning.
By modeling scenarios as spatiotemporal graphs and using a self-supervised graph autoencoder, \method injects a strong relational inductive bias. 
Evaluations show that our learned representations are highly reusable and data-efficient, and can encode complex social semantics without task supervision.
Ultimately, \method proves that structural priors can offset the scarcity of grounded HRI data, paving the way for general-purpose social foundation models in robotics.


% The acknowledgments are automatically included only in the final and preprint versions of the paper.
\acknowledgments{}

\clearpage

%===============================================================================

% no \bibliographystyle is required, since the corl style is automatically used.
\bibliography{refs}  % .bib

\end{document}
